% RESEARCH_DRAFT_SHA256: 1297a19ac64c4c242349494ce683c7e04166f0620760bde01d07e5e3d2d6c030
% Layout and source-reference wrappers expanded for the web renderer.
\documentclass[11pt]{article}
\usepackage[margin=25mm]{geometry}
\usepackage{amsthm}
\usepackage{fontspec,unicode-math}
\setmainfont{Latin Modern Roman}
\setmathfont{Latin Modern Math}
\usepackage{amsmath,amssymb,mathtools,xurl,hyperref}
\hypersetup{colorlinks=true,urlcolor=blue,linkcolor=blue}
\setcounter{secnumdepth}{0}
\setlength{\parindent}{0pt}
\setlength{\parskip}{5pt}
\setlength{\emergencystretch}{3em}
\newcommand{\sref}[2]{\hyperlink{source:#1:#2}{[S#2]}}
\newcommand{\cataloguescope}[1]{\paragraph{Source record.}#1}
\newcommand{\floor}[1]{\left\lfloor#1\right\rfloor}
\newtheorem{proposition}{Proposition}
\newtheorem{lemma}[proposition]{Lemma}
\begin{document}
\paragraph{Research manuscript.} OpenAI \texttt{gpt-6-astra}, reasoning \texttt{ultra}; 11 October 2026.
\paragraph{Preserved source body.} The complete catalogue mathematical text and sources follow with layout and source links expanded, including source qualifications and its historical preparation notes. New research begins at the section headed Research attempt.

% UnsolvedMath ID 30000195; source catalogue code OWR-792-008
\section{Interpolation-Point Selection in Rational Krylov Reduction}\label{30000195}
{\small UnsolvedMath ID 30000195 · Computer Science}
\subsection{Definitions and mathematical statement}

\paragraph{Shared notation.}Standard number systems are $\mathbb N,\mathbb Z,\mathbb Q,\mathbb R,\mathbb C$. Specialized notation and index ranges follow the cited source and the definitions below.


For a linear system with state matrix $A$, rational Krylov spaces use shifted resolvents such as $(A-\sigma I)^{-1}b$. Projection reduces the state to a subspace; interpolation points are the shifts or transfer-function frequencies used to construct that subspace.
\paragraph{Record scope and qualifications.}
The catalogue title is retained as metadata. Named certificates, contracts, bounds or frameworks in that title are not asserted as results of the original source extract.
The desired properties, interpolation-point constraints, and response-error criterion are not specified in the source question.

Source-record limitation: Desired reduced-system properties and the approximation criterion are unspecified; this is a qualitative design prompt.
\paragraph{Statement recorded in the supplied source.}
Given a large-scale linear system and a reduced order, how should interpolation points be chosen in projection-based rational Krylov model reduction so that the reduced system has desired properties? \href{https://doi.org/10.4171/owr/2005/11}{[S2]}
\paragraph{Source status.}
The supplied catalogue labels this source record open; that metadata was checked against the supplied snapshot on 2026-10-11. The label is preserved for provenance and is not an independent verification of present open status or mathematical correctness. The catalogue literature assessment is dated 2026-08-21.
\subsection{Short English statement}
Choose reduction interpolation points to obtain the desired response properties.
\subsection{Sources}
\begin{itemize}
\item[S1] Ulam AI and the UnsolvedMath Contributors, UnsolvedMath: A Curated Collection of Open Mathematics Problems, record 30000195 (OWR-792-008). Catalogue curation and metadata: CC BY 4.0. The source extract is retained; definitions, English summaries and qualifications are editorial additions.
\url{https://huggingface.co/datasets/ulamai/UnsolvedMath}.
\item[S2] Regelungstheorie (2005), Oberwolfach Reports, DOI 10.4171/owr/2005/11, pp. 565--645
\url{https://doi.org/10.4171/owr/2005/11}.
\item[S3] projection and interpolation context
\url{https://doi.org/10.1515/9783110199741.75}.
\item Source-specific Oberwolfach Reports licences, including CC BY-SA 4.0 where applicable, are retained; see THIRD\_PARTY\_NOTICES.md.
\end{itemize}
\paragraph{Preparation and review.}
The original source and UnsolvedMath contributors retain attribution for the source extract. Codex (GPT-6; exact serving revision not exposed) assisted with definitions, English summaries, scope notes and formatting. No new human mathematical authorship or independently verified solution is asserted. Permanent source ID 30000195 and code OWR-792-008 checked against the complete supplied registry and current repository. Distinct upstream records, including repeated questions and broad or damaged extracts, are retained for complete coverage of the source CS category. A repeated record is not a new independent conjecture. Complete mathematical-equivalence and current-literature review has not been performed. Preparation and final review follow contribution prompts C and D, with source gaps disclosed. Structural and rendering checks do not establish a new theorem or open status.
\label{end:30000195}

\section{Research attempt}
\subsection{Literature, source context and scope}
The full original catalogue body and source qualifications are preserved above. The relevant primary workshop passages were read in the cached publisher report. Historical claims and subsequent attributed results are distinguished from the mathematical work below. A restricted example, standard lemma or known theorem is not counted as new progress on the entire original target. Missing assumptions and repeated or rhetorical extracts are not silently turned into different solved problems.

\subsection{A two-point rational Krylov interpolation that loses stability}
The source leaves the desired properties unspecified and discusses projection spaces built from resolvents. We select one concrete desired property, asymptotic stability, and test whether interpolation at positive real points guarantees it. Consider
\[
A=\begin{pmatrix}-1&4\\0&-1\end{pmatrix},\quad
B=\begin{pmatrix}-1\\3\end{pmatrix},\quad C=(1,0),\quad D=0,
\quad V=\begin{pmatrix}1\\1\end{pmatrix},\quad W^T=(1/2,1/2).
\]
The right interpolation point is 2 and the left one is 3.
\begin{proposition}
These vectors form the report's one-dimensional two-sided rational Krylov projection. The full system is asymptotically stable; its reduced system is unstable, although the transfer functions agree at both chosen interpolation points.
\end{proposition}
\begin{proof}
Direct inversion gives $(2I-A)^{-1}B=(1,1)^T=V$ and $C(3I-A)^{-1}=(1/4,1/4)=\overline W^T$. Since $\overline W^TV=1/2$, the normalization in the source gives $W^T=2\overline W^T$, and $W^TV=1$. Projection yields
$\widehat A=W^TAV=1$, $\widehat B=W^TB=1$, and $\widehat C=CV=1$.
The full state matrix has both eigenvalues $-1$ and $e^{At}=e^{-t}\begin{pmatrix}1&4t\\0&1\end{pmatrix}$, which tends to zero. Its transfer function is
\[
H(s)=C(sI-A)^{-1}B=\frac{11-s}{(s+1)^2},
\qquad \widehat H(s)=\frac1{s-1}.
\]
Both equal 1 at $s=2$ and $1/2$ at $s=3$. The reduced zero-input state evolves as $e^t\widehat x(0)$ and is unstable. Neither chosen point is a pole of either transfer function, so the interpolation equalities are well-defined.
\end{proof}
The attempted inference that good interpolation-point placement alone automatically preserves stability therefore fails for this permitted construction. The original controllability and observability assumptions do not rescue this example: $[B,AB]$ has determinant $-36$, and the observability matrix with rows $C,CA$ has determinant 4, both nonzero. Additional restrictions on the projection, the original system or the selected objective are needed.
\subsection{Verification and remaining gaps}
All resolvent vectors, reduced matrices, determinants and interpolation values were checked exactly with rational arithmetic. No general point-selection algorithm, error criterion, passivity guarantee or globally optimal reduction was obtained. This explicit counterexample addresses only stability and does not define the source's entire unspecified set of desired properties.

\section{References and source attribution}
Primary workshop report: \url{https://doi.org/10.4171/owr/2005/11}. Relevant passages read: 2005/11 pp. 634--635, model reduction by projection and the point-selection question.

All retrievals were checked on 11 October 2026. Original question attribution remains in the preserved Sources section. External theorems are attributed and not claimed as new discoveries or fully verified merely by citation.


\paragraph{Integration review.} Codex (GPT-6; exact serving revision not exposed) reviewed the written arguments and expanded layout and source-reference wrappers. Repository metadata and the production web renderer were checked using Python 3.13 and JavaScript with MathJax 3.2.2. This model-assisted review is not independent human proof verification; no PDF compilation or formal proof checking is asserted.

\section{Completion Estimate}
\noindent\textbf{COMPLETION: 0\%}
No substantiated new progress toward the entire original target is claimed. The worked results and failed route are delimited above; a known result, illustrative model, or auxiliary obstruction is not counted as original-target completion.

\section{Final status and attribution}
\textbf{Attempt status: unresolved.} This describes this research attempt, not a new ruling on the source catalogue's status or on external literature claims.
The mathematics and this self-audit were generated by OpenAI \texttt{gpt-6-astra}, reasoning \texttt{ultra}, on 11 October 2026. Codex (GPT-6; exact serving revision not exposed) orchestrated the batch and contributed formatting/repository review. No human mathematical authorship, independent human proof checking, or formal verification is claimed. The historical preparation paragraphs in the reproduced catalogue body concern that body only.

Contribution prompt A was applied to the complete definition and supplied source record, followed by the adversarial checks of prompt B. Local reproduction files are \texttt{batch016/checks.py}, \texttt{computations.json}, and the per-ID audit JSON. They were not public links at preparation time. Numerical checks are finite and do not prove unrestricted results. No LaTeX engine was found locally; compilation was not performed by this research agent.

\end{document}
